\section{Introduction and Problem Statement}

Practical, hands-on work is the most effective way to teach offensive web
security, and Capture-the-Flag (CTF) challenges have become the standard format
for delivering it~\cite{Chapman2014picoctf,Svabensky2020what}.
In a typical course, an instructor assigns a set of vulnerable-web-application
challenges. Each student works through them and submits a ``flag'', a secret
string that appears only after a vulnerability is exploited, as proof of
completion. The format scales, students find it motivating, and it maps cleanly
onto real penetration-testing skills. That is why it is so common.

Three problems appear once this format is used with a real class, and they are the
motivation for this project. \textbf{First, every learner is different, but the
platform is not.} A beginner and an experienced student are given the same fixed
sequence of challenges at the same fixed difficulty. Weaker students are
overwhelmed while stronger ones are bored, and neither receives assignments matched
to their level~\cite{Streicher2016adaptive,Vykopal2022smart}. \textbf{Second, game
elements are added without care.} Points, badges, and especially leaderboards are
bolted on because they are expected, yet a large body of evidence shows these same
elements can reduce motivation and performance for exactly the students who most
need support~\cite{Toda2018dark,Almeida2023negative}. \textbf{Third, the instructor
is nearly blind.} Existing platforms report that a challenge was solved, and
sometimes how long it took, but they do not turn that into a usable picture of who
is falling behind, so teacher evaluation stays manual and struggling students are
found late~\cite{Svabensky2024detecting}.

This project addresses all three. We propose to design and build an
\emph{adaptive, gamified web-security training platform} for penetration-testing
education. It hosts an intentionally vulnerable web application covering the OWASP
Top~10~\cite{Owasp2021top10}, scores each attempt automatically from completion,
accuracy, and time, and wraps the whole thing in a gamification layer that is
designed around, not in spite of, the negative-effects literature. Two ideas make
the platform \emph{adaptive} rather than static. A hint-and-resource layer supports
beginners without simply handing them the answer, and a lightweight
machine-learning model learns from each student's own performance history to
predict who is likely to struggle and to recommend a next challenge at the right
level, so that assignments are customised to the individual. An instructor
dashboard makes the resulting progress data visible at a glance. The platform is
deliberately scoped as a group project: the machine-learning component is kept
simple and interpretable, and the aim is a well-engineered, evaluated teaching tool.
The rest of this synopsis surveys the literature that motivates the work, distils
the gaps we intend to cover, states the objectives, presents the proposed system
and its methodology, and sets out the expected outcomes.
